\scanpage{64}
% letterhead
% @ upper left, printed emblem | drawing
\figure{\begin{tikzpicture}[x=.19cm,y=.19cm,line width=1.6pt,line cap=butt,line join=miter]
\draw (0,0)--(0,6) (6,0)--(6,6);
\draw (1,6)--(1,3.5)--(5,3.5)--(5,6) (2,6)--(2,4.5)--(4,4.5)--(4,6);
\draw (1,0)--(1,2.5)--(5,2.5)--(5,0) (2,0)--(2,1.5)--(4,1.5)--(4,0);
\end{tikzpicture}}

% @ upper left, sender | letter part
CAISSE\\
DE RETRAITE\\
ET DE\\
PRÉVOYANCE\\
HAUSSMANN\\
Institution AGIRC n\textsuperscript{o} 36 agréée par\\
arrêté du Ministère du Travail du 10-5-48

% @ upper right, title and date | letter part
\textbf{INFORMATION AUX RETRAITES}\\
(Mars 1991)

% @ centre, salutation | letter part
Madame, Monsieur,\\

% @ centre and left column, paragraph 1 | text
\margin{La nouvelle valeur du point}
Nous avons l'honneur de vous informer que le Conseil d'administration de l'AGIRC a porté la valeur annuelle du point à 2,179 F au 1er Janvier 1991 en adoptant une augmentation de 4,11 \% par rapport à la valeur moyenne du point pour l'année 1990.

% @ centre, paragraph 2 | text
Cette nouvelle valeur tient compte d'un recalage sur l'évolution des prix observée au cours des trois dernières années et, en particulier, de l'accélération constatée dès l'été dernier et d'une augmentation moyenne des prix de 1991.

% @ centre, paragraph 3 | text
Cette revalorisation est faite, en principe, pour toute l'année 1991.

% @ lower middle, paragraph 4 | text
L'allocation qui vous sera payée le 2 Avril 1991 au titre du 1er trimestre 1991 tiendra compte de cette nouvelle valeur de point.

% @ lower third, paragraph 5 and left column | text
\margin{La CSG}
Nous vous informons également que la Loi de finances pour 1991 a instauré une contribution sociale généralisée (C.S.G) au profit de la Sécurité sociale.

% @ lower third, paragraph 6 | text
Celle-ci doit être retenue à la source, sur les arrérages versés \textbf{à partir du 1er Février 1991.}

% @ lower third, paragraph 7 | text
En conséquence, toutes les sommes payées par nos services depuis cette date, qu'il s'agisse d'arrérages trimestriels ou de rappels, sont soumises à une retenue supplémentaire.

% @ lower third, paragraph 8 | text
Cette retenue de 1,10 \% est assise sur le montant brut de la retraite versée \textbf{(y compris les majorations pour enfants)}, et \textbf{n'est pas déductible des revenus imposables.}

% @ foot, contact line | letter part
Bureaux et Réception : 64 bis, rue de Monceau 75378 Paris cedex 08 - Tél. : (1) 45.63.79.10\\
Nos collaborateurs pratiquant l'horaire personnalisé, la réception de nos participants est assurée du lundi au vendredi de 9 h 30 à 11 h 55 et de 13 h 25 à 16 h 30 ou sur rendez-vous.
% archive sheet number: 63
